\section{Conclusions}

The project shows that the robot movement planning problem can be modeled
compactly by separating chunking and placement decisions. The interval
representation of chunks, based on start index and length, avoids larger
membership-based encodings while still capturing the gripper constraints. The
release order is then handled through a permutation variable and local
neighborhood constraints on the plate direction.

The benchmark shows that the CP-oriented solvers are the most effective on the
tested instances. \texttt{gecode} is the most stable solver, while
\texttt{cp-sat} and \texttt{chuffed} remain competitive. In contrast,
\texttt{highs} is consistently slower on this combinatorial formulation. Larger
grippers reduce the number of chunks and generally improve runtime, while dense
layouts with many small boxes remain the hardest cases.